% draw-diagram TikZ 스타일. \input으로 불러서 씀. 원본은 수정하지 않음.
\usetikzlibrary{positioning,arrows.meta,fit,calc,backgrounds,shapes.geometric}

% 팔레트 (assets/palette.md와 같은 값)
\definecolor{dgBlue}{HTML}{BFDFFB}
\definecolor{dgOrange}{HTML}{FBD2B4}
\definecolor{dgMint}{HTML}{C7F5EC}
\definecolor{dgPurple}{HTML}{DCD3F9}
\definecolor{dgYellow}{HTML}{FBEBB4}
\definecolor{dgGreen}{HTML}{D3EDC8}
\definecolor{dgPink}{HTML}{F9D1F0}
\definecolor{dgGray}{HTML}{D9D9D9}
\definecolor{dgPanel}{HTML}{F2F2F2}
\definecolor{dgLine}{HTML}{404040}

\tikzset{
  every picture/.append style={line width=0.6pt, draw=dgLine, text=dgLine,
    font=\sffamily},
  % 모듈: 채운 둥근 사각형, 테두리 없음. 색은 fill=dgXxx로 지정
  module/.style={draw=none, rounded corners=3pt, minimum height=7mm,
    minimum width=24mm, align=center, font=\sffamily\small},
  % 데이터: 흰 바탕 얇은 테두리
  data/.style={draw=dgLine, fill=white, rounded corners=2pt, line width=0.4pt,
    minimum height=7mm, inner sep=3pt, align=center, font=\sffamily\small},
  % 스칼라·변수: 원
  scalar/.style={circle, draw=dgLine, line width=0.4pt, fill=dgPink,
    minimum size=6mm, inner sep=0pt, font=\sffamily\small},
  % 변수 칩: 작은 사각형 + 기호
  chip/.style={draw=none, rounded corners=2pt, minimum size=6mm,
    inner sep=1pt, font=\sffamily\small},
  % 연산 글리프: 원 안의 기호 (+, \times, \odot 등)
  op/.style={circle, draw=dgLine, line width=0.4pt, fill=white,
    minimum size=4.5mm, inner sep=0pt, font=\sffamily\footnotesize},
  % 차원 축소 층: 사다리꼴. 좁아지는 쪽이 출력 (shape border rotate로 방향 조정)
  reduce/.style={trapezium, trapezium angle=75, shape border rotate=270,
    draw=dgLine, line width=0.4pt, fill=dgGray, minimum height=5mm,
    minimum width=12mm, inner sep=1pt, font=\sffamily\footnotesize},
  % 분기: 마름모
  branch/.style={diamond, draw=none, fill=dgLine, minimum size=3.5mm, inner sep=0pt},
  % 컨테이너: fit=(...)과 함께 background 층에
  container/.style={draw=none, fill=dgPanel, rounded corners=6pt, inner sep=3mm},
  % 반복 경계: 점선 컨테이너
  repeatbox/.style={draw=dgLine!50, dashed, fill=dgPanel, rounded corners=6pt,
    inner sep=3mm},
  % 선
  flow/.style={-{Stealth[length=2mm]}, draw=dgLine},
  ref/.style={-{Stealth[length=2mm]}, dashed},          % 참조: color=출발 색 진하게
  zoom/.style={draw=dgLine!45, dashed, line width=0.4pt}, % 확대선
  skip/.style={densely dotted, draw=dgLine},              % 생략
  % 입출력 아래 두 줄: 기호 줄, 이름 줄
  symlabel/.style={font=\normalsize, inner sep=1pt},
  namelabel/.style={font=\sffamily\small\bfseries, inner sep=1pt},
  panellabel/.style={font=\sffamily\small\bfseries, anchor=north west, inner sep=0pt},
}

% 흰 바탕: 이미지로 내보낼 때(Markdown 등) tikzpicture 끝에서 \dgCanvas 호출.
% background 층보다 아래(canvas 층)에 깔아서 컨테이너를 가리지 않음. 다크 모드에서도 보이게 함.
\pgfdeclarelayer{canvas}
\pgfsetlayers{canvas,background,main}
\newcommand{\dgCanvas}[1][3mm]{%
  \begin{pgfonlayer}{canvas}
    \fill[white] ([shift={(-#1,-#1)}]current bounding box.south west)
      rectangle ([shift={(#1,#1)}]current bounding box.north east);
  \end{pgfonlayer}}
